\section{Comparison along square-step progressions}
\label{sec:scale-transfer}

We next compare the product-space functional with the same functional on
shorter intervals.  Fixing a few prime coordinates makes its low Fourier
coefficients exact averages on a residue class.  Splitting that class into
progressions with square step preserves the forbidden integer differences
and permits a change of scale.  The comparison proved here uses no induction
hypothesis. Restriction to square-step progressions is already a key
feature of quantitative square-difference arguments; see the Main Lemma
of~\cite{PSS1988} and~\cite[Lemma~6.2]{GreenSawhney2025}. Here it is
used to compare the signed multilinear functional itself.

Keep the constants \(\gamma=1/64\) and \(\beta=10^{-5}\) fixed as above.
Enlarge the fixed prime threshold \(P\), if necessary, so that the tuple
measures are available for every prime \(p\ge P\) and
\begin{equation}
 \prod_{\substack{p\ge P\\p\ {\rm prime}}}(1+p^{-3})-1\le\frac14,
 \qquad
 |\mathfrak g(\xi)|\le q(\xi)^{-1/3}
 \label{eq:prime-threshold-conditions}
\end{equation}
whenever \(q(\xi)\) is squarefree and all its prime factors are at least
\(P\).  Such a choice is possible: the product converges, while the prime
Gauss estimate used in Proposition~\ref{prop:square-kernel} gives
\[
 \left|\frac1{p-1}\sum_{x\in\F_p^\times}e(ax^2/p)\right|
 \le\frac{\sqrt p+1}{p-1}
 =\frac1{\sqrt p-1}\le p^{-1/3}
\]
for all sufficiently large \(p\) and all \(a\ne0\) in \(\F_p\).
Chinese remaindering then gives the second inequality in
\eqref{eq:prime-threshold-conditions}, including \(q(\xi)=1\).
Define the fixed integer
\begin{equation}
 M_0=8\prod_{\substack{3\le p<P\\p\ {\rm prime}}}p.
 \label{eq:fixed-small-modulus}
\end{equation}
In particular \(M_0\ge8\).

For every positive integer \(N\), put
\begin{equation}
 \begin{gathered}
 Q=N^\beta,\qquad H=N^{\beta/2},\qquad
 \mathcal P_N=\{p\ {\rm prime}:P\le p\le Q\},\\
 \tau=\frac{\beta\gamma}{1000(d+1)},\qquad
 \sigma_0=\frac{\gamma\tau}{4}.
 \end{gathered}
 \label{eq:scale-parameters}
\end{equation}
For a square-difference-free set \(A\subseteq[N]\), the quantity that will
eventually be estimated by induction is
\begin{equation}
 Y(N,A)=
 Z_{\mathcal P_N}\bigl(\mathcal L_{[N],\mathcal P_N,Q}1_A\bigr).
 \label{eq:induction-functional}
\end{equation}
It is nonnegative by Proposition~\ref{prop:signed-holder}.  We also set
\begin{equation}
 f_s(n)=M_0\,1_{\{n\equiv s\pmod{M_0}\}}1_A(n)
 \quad(s\in\Z/M_0\Z),
 \qquad
 1_A=\E_{s\bmod M_0}f_s.
 \label{eq:small-residue-input}
\end{equation}
The factor \(M_0\) in this definition will be retained in the estimates.
All constants below may depend on \(P,d,\beta,\gamma\) and the earlier
fixed construction parameters.

\begin{lemma}[Exact coefficients on a fiber]
\label{lem:exact-fiber}
Let \(N\ge1\), let \(A\subseteq[N]\), and define \(f_s\) by
\eqref{eq:small-residue-input}.  Let \(I\subseteq[N]\) be a nonempty
interval, let \(B\subseteq\mathcal P_N\), and put
\[
 q_B=\prod_{p\in B}p,\qquad
 \mathcal S=\mathcal P_N\setminus B,\qquad D=M_0q_B.
\]
Fix \(s\bmod M_0\) and \(z_B\in X_B\), and let \(a\bmod D\) be their
combined residue class.  If
\[
 g_I=(\mathcal L_{I,\mathcal P_N,Q}f_s)(z_B,\cdot),
\]
then every character \(\xi\) on \(X_{\mathcal S}\) satisfying
\(q(\xi)q_B\le Q\) has coefficient
\begin{equation}
 \widehat g_I(\xi)
 =\frac D{|I|}
   \sum_{\substack{n\in I\\n\equiv a\pmod D}}
       1_A(n)e(-n\xi).
 \label{eq:interval-fiber-coefficient}
\end{equation}
Here \(q_\varnothing=1\), and the hat denotes the normalized Fourier
coefficient on the remaining product space.
\end{lemma}

\begin{proof}
Write \(\chi_\xi(z)=e(z\xi)\) for a product-space character.
Each original character has a unique decomposition \(\xi+\eta\), where
\(\xi\) is supported on \(\mathcal S\) and \(\eta\) on \(B\).  Their
denominators are coprime, so
\(q(\xi+\eta)=q(\xi)q(\eta)\).
The hypothesis ensures that every character \(\eta\) on \(X_B\) is
retained in the lift.  Orthogonality over this entire character group gives
\begin{align*}
 \widehat g_I(\xi)
 &=\sum_{\eta\ {\rm on}\ X_B}
     \frac1{|I|}\sum_{n\in I}
       f_s(n)e(-n(\xi+\eta))\chi_\eta(z_B)\\
 &=\frac{q_B}{|I|}
     \sum_{\substack{n\in I\\
                     n\equiv z_p\pmod p\ (p\in B)}}
       f_s(n)e(-n\xi).
\end{align*}
Substituting \eqref{eq:small-residue-input} proves
\eqref{eq:interval-fiber-coefficient}.
\end{proof}

\begin{proposition}[Comparison with a shorter interval]
\label{prop:scale-transfer}
There are constants \(C_{\rm tr}<\infty\) and \(N_{\rm tr}\) with the
following property.  Let \(N\ge N_{\rm tr}\), let \(A\subseteq[N]\) be
square-difference-free, and let \(B\subseteq\mathcal P_N\) satisfy
\(q_B\le N^\tau\).  Put
\[
 \mathcal S=\mathcal P_N\setminus B,\qquad
 D=M_0q_B,\qquad N'=\lceil N/D^2\rceil,\qquad
 \lambda=\frac{D^2N'}N.
\]
For any \(s\bmod M_0\) and \(z_B\in X_B\), define
\[
 g=(\mathcal L_{[N],\mathcal P_N,Q}f_s)(z_B,\cdot).
\]
Then \(1\le N'<N\), and
\begin{equation}
 Z_{\mathcal S}(g)
 \le \lambda^d
       \max_{\substack{A'\subseteq[N']\\
                       A'\ {\rm square\text{-}difference\text{-}free}}}
                Y(N',A')
       +C_{\rm tr}N^{-\sigma_0}.
 \label{eq:scale-transfer-diagonal}
\end{equation}
Moreover, for arbitrary choices \(s_v\bmod M_0\) and
\(z_{B,v}\in X_B\), one for every \(v\in V\), the functions
\[
 g_v=(\mathcal L_{[N],\mathcal P_N,Q}f_{s_v})(z_{B,v},\cdot)
\]
satisfy
\begin{equation}
 |\mathcal Z_{\mathcal S}((g_v)_{v\in V})|
 \le \lambda^d
       \max_{\substack{A'\subseteq[N']\\
                       A'\ {\rm square\text{-}difference\text{-}free}}}
                Y(N',A')
       +C_{\rm tr}N^{-\sigma_0}.
 \label{eq:scale-transfer-mixed}
\end{equation}
All constants are uniform in \(N,A,B,s,z_B\) and the mixed
specifications.  Finally,
\begin{equation}
 1\le\lambda<1+\frac{D^2}{N},\qquad
 \lambda^d=1+O_d(D^2/N)
           =1+O_{d,M_0}(N^{-1+2\tau}).
 \label{eq:scale-transfer-factor}
\end{equation}
\end{proposition}

\begin{proof}
We first record the parameter inequalities that make both lifts
available.  Since \(\tau<\beta/2\),
\[
 Hq_B\le N^{\beta/2+\tau}\le Q.
\]
Also
\begin{equation}
 N'\ge\frac{N}{D^2}\ge M_0^{-2}N^{1-2\tau}.
 \label{eq:shorter-scale-lower-bound}
\end{equation}
Because \(1/2-2\tau>0\), this lower bound is at least \(N^{1/2}\)
for sufficiently large \(N\).  As \(D\ge8\), we have \(N'<N\) for
every integer \(N\ge2\).  Hence, with \(Q'=(N')^\beta\),
\begin{equation}
 H\le Q'\le Q,\qquad
 N\ge10Q^6,\qquad N'\ge10(Q')^6
 \label{eq:scale-transfer-cutoffs}
\end{equation}
once \(N\) is sufficiently large.  The length conditions follow from
\(6\beta<1\) and the uniform growth of \(N'\) in
\eqref{eq:shorter-scale-lower-bound}.

\smallskip
\noindent\emph{Removing large remaining conductors.}
Split \(g=g_{\le H}+g_{>H}\) according to the denominator on the
remaining prime coordinates \(\mathcal S\).
Before specialization, these pieces retain exactly the denominator
families
\[
 \{U\subseteq\mathcal P_N:q_U\le Q,\ q_{U\setminus B}\le H\},
 \qquad
 \{U\subseteq\mathcal P_N:q_U\le Q,\ q_{U\setminus B}>H\}.
\]
Thus Lemma~\ref{lem:lift-moments} applies to both pieces.
For a function \(u\) on \(X_B\times X_{\mathcal S}\), uniform counting
gives the specialization bound
\[
 \|u(z_B,\cdot)\|_r
 \le q_B^{1/r}\|u\|_r
 \qquad(r\ge1);
\]
this is also the estimate in Lemma~\ref{lem:lift-specialization}.
Since \(|f_s|\le M_0\), it follows that, for
\(g_*\in\{g,g_{\le H},g_{>H}\}\) and every fixed even \(r\),
\begin{equation}
 \|g_*\|_r\le M_0q_B^{1/r}N^{o(1)}.
 \label{eq:specialized-lift-moments}
\end{equation}
All subpower bounds in this proof are uniform over the indicated
sets and fibers.

Telescope the \(d\) inputs in the diagonal integral from \(g\) to
\(g_{\le H}\).  Each difference has one input \(g_{>H}\).
Lemma~\ref{lem:conductor-decay} and
\eqref{eq:specialized-lift-moments} imply, for every fixed
\(\varepsilon>0\),
\begin{equation}
 |Z_{\mathcal S}(g)-Z_{\mathcal S}(g_{\le H})|
 \le C_\varepsilon H^{-\gamma}q_BN^\varepsilon.
 \label{eq:scale-transfer-first-error}
\end{equation}
Indeed the specialization powers in that lemma add to
\[
 \frac12+\frac{d-1}{2(d-1)}=1.
\]
The factors \(d\) and \(M_0^d\) are fixed and are included in
\(C_\varepsilon\).  Applying the moment lemma with sufficiently small
individual exponents makes the aggregate moment loss at most
\(N^\varepsilon\).

\smallskip
\noindent\emph{Splitting the exact fiber into square-step progressions.}
Let \(a\bmod D\) combine \(s\) and \(z_B\).  By
Lemma~\ref{lem:exact-fiber}, for every remaining \(\xi\) with
\(q(\xi)\le H\),
\begin{equation}
 \widehat g(\xi)=\frac D N
   \sum_{\substack{n\le N\\n\equiv a\pmod D}}1_A(n)e(-n\xi).
 \label{eq:scale-transfer-fiber}
\end{equation}
Let
\[
 \begin{gathered}
 J_a=\{j\in[D^2]:j\equiv a\pmod D\},\qquad
 c_j=j-D^2,\\
 A'_j=\{n'\in[N']:D^2n'+c_j\in A\}.
 \end{gathered}
\]
There are exactly \(D\) elements of \(J_a\).
The map \((n',j)\mapsto D^2(n'-1)+j\) bijects
\([N']\times[D^2]\) with \([D^2N']\).
Restricting \(j\) to \(J_a\) selects exactly the residue class
\(a\bmod D\).  The definition of \(A'_j\) assigns zero to all padded
points greater than \(N\), so this decomposition incurs no additive
error.

Each \(A'_j\) is square-difference-free: a difference \(m^2>0\)
between its elements would give the difference
\(D^2m^2=(Dm)^2\) between the corresponding elements of \(A\).
This argument concerns actual integer squares.

For \(p\in\mathcal S\), the integer \(D\) is invertible modulo \(p\).
Define
\[
 T_j(z)_p=D^{-2}(z_p-c_j)\pmod p,\qquad
 h_j^H=\mathcal L_{[N'],\mathcal S,H}1_{A'_j}.
\]
Each \(T_j\) is a translation followed by multiplication by the
nonzero square \((D^{-1})^2\).
Multiplication of a remaining frequency by \(D^2\) preserves its
exact denominator, so changing variables \(\eta=D^2\xi\) in the
Fourier expansion of \eqref{eq:scale-transfer-fiber} gives the exact
identity
\begin{equation}
 g_{\le H}(z)
 =\frac{\lambda}{D}\sum_{j\in J_a}h_j^H(T_j(z)).
 \label{eq:scale-transfer-affine-average}
\end{equation}
For clarity, the coefficient of \(e(z\xi)\) on the right is
\[
 \frac{DN'}N\sum_{j\in J_a}e(-c_j\xi)
 \left(\frac1{N'}\sum_{n'\le N'}
       1_{A'_j}(n')e(-D^2n'\xi)\right),
\]
which is exactly \eqref{eq:scale-transfer-fiber}.

The seminorm triangle inequality and affine invariance in
Lemma~\ref{lem:functional-properties} now imply
\[
 Z_{\mathcal S}(g_{\le H})^{1/d}
 \le \frac{\lambda}{D}
        \sum_{j\in J_a}Z_{\mathcal S}(h_j^H)^{1/d}
 \le \lambda\max_{j\in J_a}Z_{\mathcal S}(h_j^H)^{1/d}.
\]
Consequently,
\begin{equation}
 Z_{\mathcal S}(g_{\le H})
 \le\lambda^d\max_{j\in J_a}Z_{\mathcal S}(h_j^H).
 \label{eq:scale-transfer-progressions}
\end{equation}
The averaging by \(D\) is essential: the number of progressions
causes no further loss.

\smallskip
\noindent\emph{Recovering the smaller-scale functional.}
Set
\[
 h_j^{Q'}=\mathcal L_{[N'],\mathcal S,Q'}1_{A'_j}.
\]
The difference \(h_j^{Q'}-h_j^H\) has conductor greater than \(H\).
By \eqref{eq:scale-transfer-cutoffs}, the moment lemma applies at
scale \(N'\) to these lifts and their pieces.  A second telescoping
application of Lemma~\ref{lem:conductor-decay}, this time without
specialization, gives
\begin{equation}
 |Z_{\mathcal S}(h_j^H)-Z_{\mathcal S}(h_j^{Q'})|
 \le C_\varepsilon H^{-\gamma}(N')^\varepsilon
 \label{eq:scale-transfer-second-error}
\end{equation}
uniformly in \(j\).

The lift \(h_j^{Q'}\) depends only on coordinates in
\(\mathcal S'=\mathcal P_{N'}\setminus B\), because every prime in a
denominator at most \(Q'\) is itself at most \(Q'\).
On the other hand, the full lift
\[
 u'_j=\mathcal L_{[N'],\mathcal P_{N'},Q'}1_{A'_j}
\]
becomes exactly \(h_j^{Q'}\) after uniform averaging over
\(B\cap\mathcal P_{N'}\).  Character orthogonality removes precisely
the denominators containing one of those primes.
By Lemma~\ref{lem:functional-properties}, coordinate averaging contracts
the seminorm; unused probability factors integrate to one.  Therefore
\begin{equation}
 Z_{\mathcal S}(h_j^{Q'})
 =Z_{\mathcal S'}(h_j^{Q'})
 \le Z_{\mathcal P_{N'}}(u'_j)
 =Y(N',A'_j).
 \label{eq:scale-transfer-projection}
\end{equation}
This also covers primes of \(B\) larger than \(Q'\), which are absent
from \(\mathcal P_{N'}\) already.

Combining \eqref{eq:scale-transfer-first-error},
\eqref{eq:scale-transfer-progressions},
\eqref{eq:scale-transfer-second-error}, and
\eqref{eq:scale-transfer-projection} proves
\eqref{eq:scale-transfer-diagonal} with error
\begin{equation}
 C_\varepsilon\left(
 H^{-\gamma}q_BN^\varepsilon+
 \lambda^dH^{-\gamma}(N')^\varepsilon\right).
 \label{eq:scale-transfer-explicit-error}
\end{equation}

\smallskip
\noindent\emph{Uniform error bounds and mixed inputs.}
The definition of \(N'\) gives
\(1\le\lambda<1+D^2/N\), and
\[
 D^2/N\le M_0^2N^{-1+2\tau}=o(1).
\]
Since \(d\) is fixed, these inequalities prove
\eqref{eq:scale-transfer-factor} and show that \(\lambda^d\)
is bounded for large \(N\).
Choose \(\varepsilon=\beta\gamma/8\) in
\eqref{eq:scale-transfer-explicit-error}.  Its two terms, before
the bounded factor \(\lambda^d\), are at most
\[
 N^{-3\beta\gamma/8+\tau}
 \quad\hbox{and}\quad N^{-3\beta\gamma/8},
\]
respectively.  As
\(\tau\le\beta\gamma/1000\) and
\(\sigma_0=\gamma\tau/4\le\tau/4\), they are
\(O(N^{-\sigma_0})\).  This establishes the diagonal bound with
constants independent of any inductive estimate for \(Y\).

Every \(g_v\) in the mixed statement is real and satisfies the same
diagonal bound.  Its right-hand side is nonnegative and independent
of \(v\).  Proposition~\ref{prop:signed-holder} therefore gives
\eqref{eq:scale-transfer-mixed}.
\end{proof}

In particular, the mixed estimate holds pointwise after specifying all
labels at primes in \(B\), regardless of the law used to choose those
labels.  It can therefore be integrated against any such probability
law.  When the prime laws form a product, the remaining coordinates
still carry the full laws used in Proposition~\ref{prop:scale-transfer}.
No reflection positivity is required of the law on the specified labels.
